% Generated from project outputs. Do not edit by hand.
\begin{table}[!htbp]
\centering
\begin{threeparttable}
\caption{Labor income by subsample: 2022 moments}
\label{tab:annual-2022-y1-moments}
\small
\setlength{\tabcolsep}{4pt}
\renewcommand{\arraystretch}{1.15}
\begin{tabularx}{\linewidth}{@{}Xrrrrr@{}}
\toprule
Sample & N & Mean & SD & Min & Max \\
\midrule
General & 94 & 3,232.9 & 1,595.5 & 608.3 & 7,211.2 \\
Men & 94 & 3,657.2 & 1,780.3 & 755.4 & 9,591.9 \\
Women & 94 & 2,572.6 & 1,604.2 & 251.5 & 7,268.1 \\
Urban & 94 & 3,428.8 & 1,699.5 & 298.8 & 7,602.9 \\
Rural & 94 & 3,074.6 & 2,099.0 & 473.6 & 15,700.0 \\
Indigenous & 94 & 3,077.5 & 2,272.5 & 344.8 & 17,506.3 \\
Non-Indigenous & 94 & 3,295.7 & 1,648.8 & 413.7 & 7,330.2 \\
Bottom 99 percent & 94 & 3,000.2 & 1,303.0 & 621.8 & 6,099.9 \\
Bottom 90 percent & 94 & 2,446.4 & 917.8 & 671.6 & 4,468.2 \\
Bottom 50 percent & 94 & 1,397.6 & 574.2 & 0.0 & 3,456.6 \\
\bottomrule
\end{tabularx}
\begin{tablenotes}[flushleft]
\footnotesize
\item Monthly quetzales. Unweighted distribution of survey-weighted cohort means before transition matching. N counts cohort cells, not individuals. SD is dispersion, not a standard error. Subsamples overlap. Cumulative income definitions and treatment of missing components follow the merging scripts.
\end{tablenotes}
\end{threeparttable}
\end{table}
